% SGA 3, Expose VI_A, Loop 1 English reconstruction.
% Scope: paragraphs 3.2.4 and 3.2.5 in full; subsection 3.3 is excluded.
% Authority: Polo--Gille Exp6A-13oct24.pdf, component-local pages 24--25
% (combined-reader pages 327--328; printed pages 316--317).
% OCR and the Reinhold English Markdown are comparison witnesses only.
% The Cartesian square is an exact 180-dpi crop of the controlling PDF:
% local p. 25, x=540,y=575,w=450,h=145 on a 1488x2105 render.

\medskip
\phantomsection\label{par:sga3-VIA-3.2.4}%
\noindent\textbf{3.2.4.---}
为完成在 \(F\) 和 \(G\) 在 \(A\) 上有限型时对
\hyperref[thm:sga3-VIA-3.2]{3.2} 的断言 (i), (ii), (iii) 和 (iv)(a) 的证明,
只需研究商 \(F\backslash G\).
由\SGARefLink{lem:sga3-V-6.1}{第 V 篇, 6.1},
断言\SGARefLink{thm:sga3-VIA-3.2:item:ii}{(ii)},
\SGARefLink{thm:sga3-VIA-3.2:item:iii}{(iii)}
和\SGARefLink{concl:sga3-VIA-3.2-iv}{(iv)}\SGARefLink{thm:sga3-VIA-3.2:item:a}{(a)}
在下述基变换后``成为真的'':
\[
  f:\operatorname{Spec}A'\longrightarrow\operatorname{Spec}A.
\]
因此由 EGA~IV\textsubscript{2}, 2.6.1, 2.6.2 和 2.7.1,
这些断言在基变换前也成立.
最后, 为证明 (i) 的第二个断言, 即 \(F\backslash G\)
是 \((\operatorname{pr}_2,\lambda)\) 在所有赋环空间组成的范畴中的余核,
只需参见第 V 篇, \S6(c).

\medskip
\phantomsection\label{par:sga3-VIA-3.2.5}%
\noindent\textbf{3.2.5.---}%
\footnote{\label{note:sga3-VIA-41}编者注: 此段为新增内容.}
现在仿照第 VI\textsubscript{B} 篇\SGARefLink{lem:sga3-VIB-9.2.1}{9.2.1} 的证明,
证明\SGARefLink{thm:sga3-VIA-3.2}{3.2} 的断言 (iv)(c). 记
\[
  X=F\backslash G
\]
并设 \(d\) 为以 \(\lambda\) 和 \(\operatorname{pr}_2\) 为分量的态射
\[
  F\times_A G\longrightarrow G\times_A G
\]


由\hyperref[prop:sga3-VIA-2.5.2]{2.5.2}, \(u\) 为闭浸入, 且由于
\[
  d=\sigma\circ(u\times\operatorname{id}_G),
  \qquad
  \sigma(x,y)=(xy,y),
\]
其中 \(\sigma\) 为 \(G\times_A G\) 的自同构, 态射 \(d\) 为闭浸入.
另一方面, 由 (iv)\SGARefLink{prop:sga3-VIA-2.5.2:item:a}{(a)}, 下图为笛卡儿方图:
\begin{center}
  \phantomsection\label{diag:sga3-VIA-3.2.5-cartesian-square}%
  \begin{tikzcd}[row sep=large,column sep=large]
    F\times_A G \arrow[r,"d"] \arrow[d] &
    G\times_A G \arrow[d,"p\times p"] \\
    X \arrow[r,"\Delta_X"'] &
    X\times_A X
  \end{tikzcd}
\end{center}
此外, \(p\), 从而 \(p\times p\), 忠实平坦且局部有限呈示.
因此, 由 fppf 下降, 由于 \(d\) 为闭浸入, \(\Delta_X\) 也为闭浸入;
即 \(X\) 分离.
% Independent prose review: source and translation compared paragraph by paragraph through this component; statements, hypotheses, proofs, notes and diagram captions covered; bibliography reserved for its assigned reviewer; source math and reference keys unchanged.
